\section*{Big Picture}

Attention is unevenly distributed across populations, and so is our ability to study it.
Individual differences in sustained attention and working memory predict a wide range of
consequential life outcomes, from academic achievement and occupational performance to
vulnerability to clinical conditions like attention-deficit/hyperactivity disorder
\parencite{kofler2013}. The neuroscientific tools most capable of measuring these
differences, however, are expensive, stationary, and concentrated in well-funded research
institutions. Functional magnetic resonance imaging (fMRI) scanners cost millions of
dollars to purchase and maintain, require specially shielded rooms, and are largely absent
from the low-income communities and under-resourced clinical settings where attentional
difficulties are most prevalent. As a result, the knowledge base about attention and the
brain reflects a structurally skewed sample of who can access a scanner \parencite{farah2006}.

This project asks whether that constraint can be relaxed. The central question is whether,
if a participant wears a portable brain sensor while watching a movie, the recording from
that single naturalistic session can be used to train a computational model that later
predicts whole-brain neural activity during structured attention tasks. The portable sensor
in question is functional near-infrared spectroscopy (fNIRS), a lightweight headset that
measures cortical hemodynamics by detecting changes in light absorption through the scalp.
If the approach works, it would allow researchers to approximate fMRI-quality, whole-brain
measures of cognitive engagement from a low-cost recording taken in a classroom, a clinic,
or a community center, without any scanner access.

The design draws on two datasets linked by a shared naturalistic stimulus, the 1959
Hitchcock film \textit{North by Northwest} (NNW). An ongoing in-lab fNIRS study
(N = 20 participants collected to date, target N = 40) records oxygenated hemoglobin
(HbO) time series from 42 prefrontal, lateral parietal, and temporal-occipital channels
while participants watch the film and then complete a series of attention tasks. These
tasks include a working memory n-back task at 1-back and 2-back difficulty levels and
two rounds of a Continuous Performance Task (CPT). A separate, open fMRI dataset from a
collaborating lab (N = 59 different participants) provides fully preprocessed BOLD time
series across 122 brain regions from participants who watched the same film in a scanner.
Because both groups watched the same movie, their neural responses are temporally aligned
even though they are entirely different people. This shared temporal structure is the
engine of the approach.

The computational method at the core of the project is adapted principal component
regression (aPCR), following \textcite{gao2025}. The model first reduces both the fNIRS
and fMRI recordings to their most important underlying components via PCA, capturing 90\%
of variance in approximately 25 fNIRS components and 66 fMRI components. It then uses
ordinary least-squares regression to learn the linear mapping between those component
spaces. Training uses leave-one-out cross-validation across the 15 currently analyzed
participants. In each fold, 14 fNIRS participants are paired with all 59 fMRI
participants, producing 826 training pairs concatenated to 731,836 timepoints, and the
held-out fNIRS participant's data are passed through the frozen learned mapping to
generate a predicted whole-brain fMRI time series for 122 regions over 886 seconds.
The NNW pipeline has already been run end-to-end. The model significantly predicts
22 of 122 brain regions (FDR-corrected, $q < 0.05$), with a median Pearson $r = 0.220$
across significant regions and a peak $r = 0.310$ in left dorsal attention
temporooccipital cortex. At the connectivity level, the Pearson correlation between
predicted and true inter-subject functional connectivity (ISFC) matrices is $r = 0.530$
($p = 0.001$, permutation test), which indicates that the model preserves not just local
activation patterns but the large-scale organization of brain-wide communication during
the film. The dorsal attention and visual networks show the strongest prediction (7 of 14
and 5 of 10 significant ROIs, respectively), and both are primary targets of the n-back
and CPT tasks planned for the next phase.

The project goes beyond the original \textcite{gao2025} framework by reframing the
primary question. Rather than asking whether the model can transfer across paradigms in
a technical sense, this project poses a social science question: can we detect meaningful
individual differences in who pays attention, and can we do so without a scanner? For
each cognitive task, the frozen NNW model will be applied to each participant's task
fNIRS data to generate predicted whole-brain time series. These predictions will be
evaluated for neural plausibility, specifically whether predicted activation patterns in
frontoparietal networks exceed default mode network activation during high-load
conditions, as the fMRI literature would predict. They will also be evaluated for
individual differences, specifically whether per-subject predicted activation patterns
correlate with behavioral performance including n-back accuracy, CPT hit rate, and
reaction time variability. This is where the equity argument becomes concrete. If a
brief, portable, naturalistic recording session can approximate the neural signature of
attentional capacity, then rigorous cognitive neuroscience becomes possible in the
settings and populations that have historically been excluded from it. Reaction time
variability is particularly well-suited as the individual-differences target because it
is sensitive enough to discriminate clinical and subclinical attentional profiles, and if
the predicted neural patterns track it, the approach offers a path toward accessible,
whole-brain-equivalent measures of attentional stability in schools, clinics, and
communities where MRI infrastructure simply does not exist.

The project is advised by Yuan Chang Leong at the University of Chicago, whose lab
collected the NNW fNIRS dataset and whose prior work established the aPCR framework
\parencite{gao2025}. All data, software, and high-performance computing resources are
in place, and the NNW pipeline has been validated end-to-end. The remaining work,
including preprocessing the cognitive tasks, applying the frozen model, and running the
Phase 5 analyses, is computationally tractable given the existing infrastructure.
